% !TeX root = ../main.tex

\chapter{Conclusions and Recommendations}
\label{chap:conclusions_and_recommendations}

This chapter synthesizes the empirical and architectural findings of the study into high-level insights and sets the direction for future research in legal informatics and natural language processing. Section~\ref{sec:conclusions} presents the principal conclusion regarding the primary research problem and systematically details the empirical evidence demonstrating how each specific research objective was achieved. Section~\ref{sec:future_work} establishes the operational and technical roadmap for subsequent investigations, structured into concrete methodological limitations and technical extensions. In doing so, this chapter bridges computational theory with practical municipal administration, ensuring that the contributions of this thesis are clearly articulated for academic researchers and local governance practitioners. By contextualizing the system within the operational realities of Philippine local government units, the discussion demonstrates how local artificial intelligence can assist legislative drafting without compromising legal precision.


\section{Conclusions}
\label{sec:conclusions}

Automating the consistency review of local legislation against superior national statutes addresses a pressing institutional challenge in Philippine local governance. Under the doctrine of statutory preemption established in \textit{Magtajas v. Pryce Properties} and Section~5(a) of Republic Act No.~7160 (the Local Government Code of 1991), municipal ordinances that contradict national laws are \textit{ultra vires} and legally invalid. In local legislative bodies such as the Davao City Sangguniang Panlungsod, manual ex-ante auditing places substantial cognitive strain on legal researchers due to the sheer volume of positive law, historical scanning noise, and compressed committee schedules. This study investigated whether a multi-stage architecture could automate this auditing process on resource-constrained municipal hardware. By formulating an offline, privacy-preserving pipeline, the research directly resolves the tension between neural inference demands and local government computing limits.

The principal conclusion of this study is that a cascaded coarse-to-fine architecture, decoupling broad candidate discovery from deep logical inference, provides an accurate computational solution for ex-ante legal conflict detection on standard local government workstations. By combining hybrid sparse-dense information retrieval in Stage~1 with cost-sensitive cross-encoder inference in Stage~2, the proposed system successfully identifies statutory preemption risks without requiring cloud infrastructure, recurring API subscriptions, or generative text models. The empirical evaluations conducted across the unified dual statutory knowledge base, the faculty adviser-verified ground truth benchmark, and authentic Supreme Court jurisprudence provide concrete evidence verifying the achievement of the research objectives. Furthermore, the findings demonstrate that domain-specific calibration aligns neural decision boundaries with legal risk tolerances, preventing high-stakes statutory conflicts from escaping scrutiny. Consequently, the research demonstrates that local decision-support tools can enhance legislative vetting efficiency while operating entirely within local institutional constraints.

Regarding Specific Objective~1, the study successfully identified the linguistic characteristics, structural properties, and archival noise factors that induce retrieval failure in the Davao City municipal collection. An empirical examination of 1,925 local ordinances (comprising 13,874 provisions across four legislative eras and six functional typologies) revealed marked document length asymmetry and structural dispersion. While 94.75\% of substantive municipal provisions comply with standard 512-token transformer windows, historical enactments dating prior to the year 2000 present substantial physical wear, with 68.8\% produced via manual impact typewriters exhibiting ribbon fading and carbon bleed-through. The study demonstrated that traditional optical character recognition engines induce cascading errors on statutory texts, corrupting penalty numerals and legal citations. Deploying a locally hosted 4-bit vision-language model (\texttt{Qwen2.5-VL-3B-Instruct}) resolved these layout and character artifacts, achieving 100.0\% legal marker precision and 100.0\% semantic cosine fidelity on authentic historical ordinances while ensuring total data privacy compliance under Republic Act No.~10173.

Regarding Specific Objective~2, the study successfully designed a Coarse-to-Fine computational pipeline that optimizes resource utilization for edge deployment. Decoupling candidate retrieval from cross-attention inference reduced the active search space from 27,357 total enactments (178,494 provisions) down to a compact shortlist of top-$k$ candidates ($k \le 50$), capturing 100\% of controlling statutory preemption provisions on landmark jurisprudential test cases. To eliminate the processing latency of evaluating full-length municipal documents on consumer CPUs, the study formulated three structural optimizations: procedural boilerplate filtering, dynamic token padding, and cross-section flattened mini-batching. Empirical latency profiling demonstrated that these optimizations accelerate full-document audit throughput by $5.79\times$, reducing end-to-end execution time from 73.47~seconds down to 12.68~seconds per draft ordinance on an entry-level quad-core desktop processor. These architectural adaptations ensure that thorough statutory audits can be completed during active committee deliberations without introducing operational delays.

Regarding Specific Objective~3, the study implemented and calibrated a hybrid Information Retrieval and Natural Language Inference system tailored to the Philippine legal hierarchy. In Stage~1, fusing Okapi BM25 lexical search with \texttt{all-MiniLM-L6-v2} dense embeddings via Reciprocal Rank Fusion ($k=60$) elevated candidate recall across all retrieval tiers while enforcing a statutory hierarchy priority safeguard that prioritizes governing national codes. In Stage~2, fine-tuning candidate cross-encoders eliminated the pronounced affirmative entailment bias observed in off-the-shelf foundation models, where unadapted zero-shot models suffered a 39.70\% drop in Conflict $F_1$-score. Incorporating cost-sensitive threshold calibration ($\tau^* = 0.45$) with an asymmetric false-negative penalty ratio ($C(\text{FN})/C(\text{FP}) = 5.0$) successfully aligned the neural decision boundary with the institutional cost of overlooking invalid local penalties. This calibration ensures that the classifier operates with heightened sensitivity toward potential legal violations rather than favoring statistical majority classes.

Regarding Specific Objective~4, the study evaluated the empirical reliability and statistical validity of the proposed pipeline against verified legal conflicts. When evaluated across the 350-pair faculty adviser-verified ground truth benchmark stratified across eight Macro Legal Domains and three difficulty tiers, the proposed fine-tuned \texttt{DeBERTa-v3-base-NLI} cross-encoder achieved a 5-fold cross-validation accuracy of $88.30\% \pm 3.44\%$, a Macro $F_1$-score of $0.8800 \pm 0.0371$, and a Conflict $F_1$-score of $0.9008 \pm 0.0421$. Non-parametric statistical testing confirmed that the proposed architecture significantly outperforms canonical supervised baselines (\texttt{bert-base-uncased}, McNemar's $\chi^2 = 4.27, p = 0.0388$) and maintains consistent top-rank performance across all legal domains (Friedman's $\chi_F^2 = 40.00, p < 0.001$). Furthermore, multi-tier system testing confirmed 100.0\% document-level conflict triage recall across ten full-length synthetic draft ordinances (140 substantive sections) and 100\% preemption detection accuracy on authentic Supreme Court landmark controversies, including \textit{Mosqueda v. PBGEA} and \textit{City of Davao v. First Meridian Development}. The statistical consistency observed across both synthetic and authentic test suites establishes that the model generalizes effectively across diverse regulatory subjects.

Regarding Specific Objective~5, the study developed and validated an offline, single-page web prototype interface designed for Sangguniang Panlungsod First Reading reviews. The production interface incorporates drag-and-drop document ingestion, procedural boilerplate filtering, dual-stream candidate inspection, and a side-by-side explainable AI (XAI) inspector that displays unpolluted verbatim statutory text alongside highlighted conflicting modal operators. Profiling on a representative municipal desktop (Intel Core i3-10105F, 4 threads, 8~GB RAM) confirmed an average per-pair inference latency of 120.9~milliseconds and a peak resident memory footprint of 1,416.4~MB, operating comfortably below the 8~GB municipal memory ceiling. Due to external administrative routing schedules through the Human Resource Management Office (HRMO) and the City Administrator's Office, the human usability and usefulness evaluation with the 15 Sangguniang Panlungsod legal researchers was not executed prior to the conclusion of the thesis research period. Instead, the 14-item TAM and SUS psychometric instrument and the multi-rater consensus protocol are formally codified as a post-study institutional deployment roadmap to be administered once municipal clearances are complete. Computational profiling verifies that the standalone prototype functions as an efficient, privacy-preserving screening tool ready for immediate local deployment.


\section{Future Work}
\label{sec:future_work}

While this research demonstrates the practical efficacy of cascaded semantic conflict detection for local governance, the system possesses explicit operational and architectural boundaries that provide the foundation for subsequent investigations. To guide future researchers, this section delineates the specific limitations encountered during development and outlines concrete technical extensions. By explicitly documenting these boundaries, the study maintains methodological transparency and provides clear starting points for future development. The identified limitations reflect the real-world complexities of digitizing historical municipal archives and operating within municipal computing constraints. Conversely, the proposed extensions chart a practical path toward expanding the system's reasoning capabilities, language support, and institutional integration.


\subsection{Methodological and System Limitations}
\label{sec:limitations}

The first limitation involves the physical degradation floor of historical municipal archives and mechanical typewriter artifacts. In Philippine local government units, legislative archives dating prior to the enactment of Republic Act No.~7160 in 1991 exist almost exclusively as physical carbon-copy papers stored in municipal record halls. These historical documents suffer from mechanical strike irregularities, faded typewriter ribbons, carbon transfer bleed-through from thin onionskin paper, and physical occlusions from rusted metal staples and binder hole punches. The study demonstrated that classical optical character recognition engines frequently corrupt numerical figures on these degraded pages---such as converting a \textpeso 200.00 fee into \textpeso 2,000.00 by mistaking smudges for digits, or dropping terminal zeros entirely. While deploying a locally hosted 4-bit vision-language model (\texttt{Qwen2.5-VL-3B-Instruct}) successfully resolved semi-modern scanning defects, severe acidification and overlapping bleed-through on pre-1980 records establish an empirical ceiling on autonomous vision-language extraction. Consequently, degraded archival enactments require conservative quarantine and manual verification, leaving historical corpus digitization partially dependent upon human supervisory review.

The second limitation arises from structural legislative asymmetry and the contextual token overhead of municipal administrative provisions. National statutes typically establish normative prohibitions in one or two direct sentences. In contrast, Philippine local ordinances are operational regulatory instruments: an authentic ordinance section regularly spans four to six contiguous sentences bundling the core prohibition alongside municipal office enforcement directives, Citation Ticket issuance procedures (Ordinance Violation Receipts), 72-hour No-Contest compromise settlement mechanisms, and alternative community service options. If an automated system fragments an ordinance into isolated sentences, the classifier loses the qualifying provisos and administrative exceptions, triggering false preemption alarms. Resolving this challenge required asymmetric multi-sentence span pairing and structural metadata prepending (\texttt{[Citation | Title > Chapter > Section]}). However, this metadata injection imposes an empirical overhead of approximately 22.7 tokens per candidate sequence, consuming valuable self-attention capacity within the constrained 512-token context window of lightweight cross-encoders.

The third limitation centers on the single-premise directional formulation adopted for the inference pipeline. The Stage~2 Cross-Encoder models vertical legal consistency as an isolated directional pair, matching a single national statute section ($\mathbf{s}$) against an operative draft ordinance block ($\mathbf{o}$). This structural formulation effectively captures direct vertical contradictions where an individual national statute explicitly invalidates a local provision, such as penalty escalations exceeding the \textpeso 5,000 ceiling under Section~458 of Republic Act No.~7160. However, the architecture cannot inherently detect multi-statute cumulative preemption. In authentic administrative law, a local ordinance may comply fully with each national statute when evaluated in isolation, but breach a composite regulatory ceiling formed by the intersection of multiple separate national codes, such as when municipal tax validity requires cross-referencing Section~133 of Republic Act No.~7160 against the National Internal Revenue Code (Republic Act No.~8424). Resolving cumulative preemption requires synthesizing across multi-statutory dependency networks, which exceeds the memory and relational scope of single-pair cross-attention.

The fourth limitation concerns computational latency ceilings and multi-section pre-enactment scaling constraints on municipal edge hardware. While evaluating an isolated text pair executes in approximately 120~milliseconds on an entry-level CPU, evaluating an authentic draft ordinance reveals the difference between single-pair latency and full-document audit time. A realistic draft ordinance comprises ten to fifteen substantive regulatory sections, each evaluated against a candidate shortlist of fifty retrieved national statutes, requiring 500 to 750 deep cross-encoder forward passes. On a standard municipal desktop with a quad-core processor and zero discrete GPU acceleration, unoptimized batching resulted in execution latencies exceeding two minutes (73.47~seconds), which creates noticeable friction during fast-paced committee meetings. While the algorithmic optimizations developed in this study---procedural boilerplate filtering, dynamic token padding, and cross-section flattened batching---successfully reduced audit time by $5.79\times$ to 12.68~seconds, scaling the architecture to massive omnibus codifications (such as the 434-section Davao City Revenue Code, Ordinance No.~158-05) remains computationally demanding on entry-level CPU workstations without dedicated model compression.

The fifth limitation involves the linguistic and jurisdictional boundaries of regional legislative drafting. In accordance with Executive Order No.~292 and national codification standards, the statutory archive and local ordinance records examined in this research are limited to official English-language enactments. Consequently, the automated pipeline cannot process regional Cebuano-Visayan phrases or mixed Taglish phrasing. In practical municipal governance, while the final engrossed copy of an ordinance is enacted in formal legal English, preliminary committee deliberations and early draft concept notes frequently incorporate localized vernacular expressions. Semantic ambiguity resulting from colloquial phrasing remains beyond the capability of standard English legal tokenizers, requiring manual review by committee researchers before texts can be ingested by the system. Developing multilingual capabilities tailored to Philippine regional dialects represents an essential step toward supporting the full legislative drafting lifecycle.

The final limitation stems from the institutional administrative routing required for municipal personnel engagement. Following the formal referral letter from the Sangguniang Panlungsod Secretariat to the City Administrator's Office via the Human Resource Management Office (HRMO), multi-office administrative reviews and ongoing legislative committee commitments resulted in an extended review timeline that exceeded the academic submission schedule. While the underlying computational models and legal conflict detections are fully validated against faculty adviser-adjudicated ground truth and binding Supreme Court jurisprudence, the 15-evaluator multi-annotator consensus ($\kappa$) and the 14-item TAM and SUS psychometric surveys could not be deployed within the thesis development calendar. This institutional reality highlights the practical challenges of deploying software systems in public sector environments where administrative protocols must be respected. Future operational rollouts should establish formalized inter-agency coordination mechanisms early in the project lifecycle to accommodate municipal administrative calendars.


\subsection{Concrete Technical Extensions}
\label{sec:technical_extensions}

To build upon the theoretical and empirical foundations established in this thesis, future researchers should pursue several concrete technical extensions to enhance the reasoning depth, knowledge coverage, and operational deployment of the conflict detection architecture. These extensions address the structural, computational, and institutional boundaries identified during empirical development. By advancing from isolated text classification to relational knowledge modeling, dynamic validity tracking, and fine-grained explainability, subsequent systems can capture complex multi-statute interactions that currently require human legal synthesis. Furthermore, expanding corpus acquisition and integrating the inference engine into municipal legislative databases will ground automated decision support directly in day-to-day governance. The following discussions detail each proposed extension and outline practical implementation pathways.

First, researchers should extend the Stage~2 inference engine from isolated pair classification to multi-hop statutory graph reasoning with dynamic temporal law tracking. Future work can construct a directed Philippine Statutory Knowledge Graph connecting national codes, regulatory executive orders, and local ordinances via explicit semantic edges (\textit{empowers}, \textit{amends}, \textit{repeals}, \textit{exempts}, \textit{delegates}). Relational graph-attention networks can track whether a legal provision remains in its authentic raw state, possesses amended successor versions, or has suffered express repeal under the principle of \textit{lex posterior derogat legi priori}. Furthermore, this graph structure can incorporate symbolic multi-tiered penalty constraints (such as imprisonment limits and administrative sanctions under Republic Act No.~7160) alongside localized municipal synsets mapping colloquial regulatory terms (\textit{e.g.}, \textit{trisikad} and \textit{payong-payong}) to governing statutory definitions. Integrating graph neural networks with the cross-encoder would enable the detection of cumulative preemption where an ordinance violates an interlocking statutory framework rather than an isolated statutory command.

Second, future research should expand the knowledge base beyond statutory enactments to incorporate Supreme Court jurisprudence for ex-post conflict analysis. Under the constitutional doctrine of \textit{stare decisis}, judicial interpretations rendered by the Supreme Court establish binding legal precedents that define the substantive boundaries of local legislative authority. Extending the dual-stream retrieval and inference pipeline to index landmark decisions allows the system to evaluate how judicial doctrines invalidate or uphold ordinances following enactment. This capability would enable legislative researchers to audit active city ordinances against evolving jurisprudential precedents, identifying provisions vulnerable to judicial challenge. Developing this ex-post review layer transforms the architecture into a comprehensive legislative lifecycle tool capable of evaluating both prospective drafts and enacted municipal laws.

Third, subsequent studies should prioritize continuous corpus acquisition and conduct a systemic empirical consistency audit across positive law. Because legislative archives continuously expand with each congressional and city council session, establishing automated ingestion pipelines is necessary to raise collection completeness across both national and local jurisdictions. Utilizing this expanding repository, researchers can execute an offline data science audit cross-referencing national statutes and local enactments against one another. This large-scale empirical audit can identify latent historical contradictions, redundant enactments, and unenforced implied repeals across decades of cumulative legislation. Systematizing this analytical process provides municipal record divisions with an empirical inventory of dormant conflicts requiring legislative consolidation.

Fourth, future research should implement cross-lingual and dialectal domain adaptation to support regional municipal drafting. Subsequent studies can curate a bilingual parallel corpus of committee session transcripts and local administrative notes across Cebuano-Visayan and English. Fine-tuning regional multilingual encoders (such as \texttt{xlm-roberta-base} or specialized Philippine language checkpoints) will enable the system to identify potential statutory conflicts directly within early-stage draft concepts and public consultation remarks before they are formally translated into English statutory syntax. Such linguistic expansion bridges the communication gap between grassroots public consultations and formal statutory codification. By detecting potential legal issues early in the consultative process, the system can assist legislative staff in addressing community concerns without drafting legally defective provisions.

Fifth, researchers should deepen the system's explainable artificial intelligence (XAI) architecture beyond preliminary visual cues. The current prototype provides a foundational form of interpretability by presenting side-by-side verbatim statutory text alignments and highlighting deontic modal operators such as \texttt{shall} and \texttt{prohibited}. While this baseline ensures legal fidelity, subsequent investigations should incorporate fine-grained attribution mechanisms, such as integrated gradients or token-level attention-saliency mapping. These methods can isolate and highlight the precise contradictory clauses within multi-sentence regulatory blocks rather than relying on generic lexical markers. Providing targeted attribution rationales will further reduce cognitive strain on legal researchers during high-speed committee review.

Sixth, subsequent modeling experiments should incorporate cost-sensitive loss formulations directly into the cross-encoder training objective. While this study successfully addressed affirmative entailment bias through post-hoc decision threshold calibration ($\tau^* = 0.45$), the underlying model parameters were optimized using standard cross-entropy loss. Future work can implement asymmetric cost-weighted cross-entropy or focal loss during backpropagation, penalizing false negatives more heavily than false positives throughout the optimization process. Directly embedding the institutional cost of overlooking preemption contradictions into the loss function will encourage the model to learn conservative decision boundaries natively. Comparing natively cost-sensitive models against post-hoc calibrated baselines will establish whether direct loss weighting improves generalization across ambiguous regulatory domains.

Seventh, future research should explore automated benchmark synthesis to enhance ground truth validation and stress-testing. Using the three-tier difficulty taxonomy formulated in this thesis (Direct, Latent, and Distractor pairs), subsequent studies can design programmatic perturbation pipelines that systematically alter authentic statutory provisions. By programmatically swapping deontic modal verbs, inflating fine amounts beyond statutory ceilings, and inserting subtle jurisdictional exceptions, researchers can generate extensive sets of synthetic test pairs with verifiable ground truth labels. This automated framework would enable continuous integration testing for legal language models without requiring months of intensive manual annotation. Evaluating models against synthetically hardened adversarial distributions will provide deeper insights into boundary vulnerabilities.

Eighth, subsequent research should expand the system beyond diagnostic conflict detection to prescriptive legislative drafting assistance. Future studies can integrate a locally deployed generative language model (such as a 4-bit quantized Llama-3-8B or Qwen-2.5-7B) constrained by strict legal drafting schemas. When a preemption contradiction is detected, the generative module can automatically draft compliant amendatory clauses or explicit repealing provisions conforming to the Department of the Interior and Local Government (DILG) drafting manual, suggesting calibrated penalty amounts that fall strictly within statutory ceilings under Section~458 of Republic Act No.~7160. To maintain legal accuracy, all generated amendments must be strictly grounded in retrieved national statutes, with explicit citation trails linking proposed wording to authorizing statutory clauses. This extension would transform the system from a passive diagnostic filter into an active legislative drafting assistant.

Ninth, researchers should investigate staged inference optimizations and model compression to maximize edge CPU throughput without compromising audit thoroughness. While current optimizations reduce audit latency to 12.68~seconds per draft, scaling to extensive omnibus codes requires further runtime reductions. Future implementations should experiment with staged inference architectures, utilizing lightweight preliminary classifiers to rapidly filter obvious non-conflicting candidate pairs while reserving full cross-encoder attention for borderline cases. Combining this staged approach with 8-bit integer (INT8) quantization and ONNX Runtime execution graphs can drive execution speeds low enough to support real-time typing-time auditing inside standard word processors such as Microsoft Word and LibreOffice Writer. Seamless word-processor integration embeds automated verification directly into daily legislative drafting routines without requiring external application switching.

Tenth, future development should pursue direct technical integration with the Davao City Legislative Information System of Sangguniang Panlungsod (LISSP). Currently operating primarily as an archival database for finalized enactments, the LISSP infrastructure can be enhanced by embedding the coarse-to-fine conflict detection pipeline into its core administrative software. Such an integration would transform the platform into an active digital legislative workbench that tracks ordinances across First, Second, and Third Readings, alerting committee staff to potential preemption risks whenever amendatory drafts are uploaded. Connecting the automated engine directly to official municipal repositories ensures that legislative vetting becomes a standard digital gateway in local governance. This operational integration bridges the gap between standalone academic prototypes and institutional municipal software ecosystems.

Eleventh, subsequent investigations should develop an automated legislative lineage engine to construct consolidated municipal enactments. In Philippine local governance, major regulatory frameworks are rarely enacted as self-contained static codes; rather, an original historical enactment is frequently modified by dozens of piecemeal amendatory ordinances passed across several decades. Consequently, evaluating a solitary historical document in isolation fails to capture its contemporary legal effect. An automated lineage compiler can parse amendatory directives (\textit{e.g.}, ``Section 2 of Ordinance No. 158-05 is hereby amended to read as follows''), track the chain of legislative modifications chronologically, and sequentially apply text substitutions to assemble a live, consolidated ``as-amended'' text. Evaluating this reconstructed consolidated draft against national statutes ensures that conflict detection operates on the authentic operative state of municipal law rather than obsolete historical fragments.

Twelfth, future empirical benchmark studies should recalibrate human evaluation workload distributions to mitigate cognitive fatigue among domain experts. In the current study design, evaluating a 350-pair ground truth with threefold inter-annotator redundancy required allocating 70 dense legal pairs to each of the fifteen legal researchers. Authentic legal texts contain high syntactic complexity and multi-sentence regulatory provisos, meaning that large batch allocations place heavy cognitive strain on practicing legal officers who must balance evaluation with daily council duties. Future research should investigate fractional-block experimental designs and adaptive micro-tasking, empirically recalculating the optimal number of items per annotator (\textit{e.g.}, fifteen to twenty high-leverage pairs per session) to preserve annotator focus and maximize inter-rater reliability ($\kappa$). In parallel, the user interface should implement confidence-gated severity tiers to suppress low-priority informational flags during high-tempo committee deliberations, ensuring that decision support remains actionable without overwhelming legislative staff.

Finally, upon completion of the institutional clearance process with the Davao City Human Resource Management Office and City Administrator, researchers should execute a longitudinal 6-month field study within the Sangguniang Panlungsod committee review workflow. To ensure successful execution, future academic investigations must incorporate comprehensive operational foresight into their project timelines, mapping municipal administrative clearance hierarchies well in advance of field deployment. Navigating multi-office routing requires scheduling generous temporal buffers in research calendars to accommodate municipal administrative pacing and ongoing legislative committee schedules. Furthermore, research teams should employ an empathic and persistent liaison strategy that balances formal digital transmissions with proactive in-person office coordination, recognizing that municipal legal personnel face demanding public service commitments. Deploying the prototype dashboard during the Second Reading preparation phase across standing committees will enable long-term measurement of actual workflow adoption, capturing empirical distributions of perceived usefulness and amendment quality over time.
